% !TeX program = xelatex
% 原册第91页第8题：保留正方形轮廓，把所用对称轴画为延伸的虚线。
\documentclass{article}
\usepackage{geometry,fontspec,tikz,xcolor}
\geometry{paperwidth=595bp,paperheight=959.62bp,margin=0bp}
\setmainfont{FandolKaiSC-Regular.ttf}[Path=fonts/]
\newfontfamily\kai{FandolKaiSC-Regular.ttf}[Path=fonts/]
\definecolor{redpen}{rgb}{0.85,0.04,0.07}
\definecolor{gridink}{rgb}{.32,.34,.45}
\pagestyle{empty}
\begin{document}
\null
\begin{tikzpicture}[remember picture,overlay]
  \begin{scope}[shift={(current page.north west)},x=1bp,y=-1bp]
    \fill[white] (324.5,659) rectangle (531,781);
    \begin{scope}[shift={(325.667,660)}]
      \draw[gridink,line width=.35bp,step=17bp] (0,0) grid (204,119);
      \draw[gridink,line width=.55bp] (0,0) rectangle (204,119);
      \foreach \sx in {17,119}{
        \draw[redpen,line width=.9bp] (\sx,17) rectangle ++(34,34);
        \draw[redpen,line width=.75bp,dash pattern=on 3bp off 2bp]
          (\sx-6,11) -- (\sx+40,57);
      }
      \node[anchor=west,fill=white,inner sep=1.4bp,text=redpen,
        font=\kai\fontsize{8.5bp}{11bp}\selectfont]
        at (10,96) {虚线表示所用的对称轴。};
    \end{scope}
  \end{scope}
\end{tikzpicture}
\end{document}
